\documentclass[11pt]{article}
\usepackage[T1]{fontenc}
\usepackage{lmodern,microtype}
\usepackage[margin=1in]{geometry}
\usepackage{amsmath,amssymb,amsthm,mathtools}
\usepackage[colorlinks=true,linkcolor=blue,citecolor=blue,urlcolor=blue]{hyperref}
\hypersetup{pdftitle={Positive lower density of primitive-root primes},
 pdfsubject={A short reduction using one large predecessor factor and a common Hecke zero-free strip}}
\newtheorem{theorem}{Theorem}[section]
\newtheorem{proposition}[theorem]{Proposition}
\newtheorem{lemma}[theorem]{Lemma}
\newtheorem{corollary}[theorem]{Corollary}
\theoremstyle{remark}
\newtheorem{remark}[theorem]{Remark}
\newcommand{\ord}{\operatorname{ord}}
\newcommand{\Li}{\operatorname{Li}}
\newcommand{\1}{\mathbf 1}
\newcommand{\supp}{\operatorname{supp}}
\newcommand{\e}{\mathrm e}
\newcommand{\B}{\mathcal B}
\newcommand{\E}{\mathbb E}
\newcommand{\Q}{\mathbb Q}
\newcommand{\Z}{\mathbb Z}
\DeclareMathOperator{\Gal}{Gal}
\DeclareMathOperator{\disc}{disc}
\title{Positive lower density of primitive-root primes}
\author{AI-assisted research draft\thanks{AI assistance: OpenAI ChatGPT was used for mathematical exploration and drafting; OpenAI Codex was used for editing and compilation.}}
\date{8 October 2026}

\hypersetup{pdfsubject={Research draft}}
\begin{document}
\maketitle

\begin{abstract}
We give a short deduction of positive lower density in Artin's
primitive root problem from a common Hecke zero-free strip and the
marked prime-predecessor estimates of OpenAI.  For every integer
$a$ that is neither $-1$ nor a square, the resulting bound is
$\#\{p\le X:\ord_p(a)=p-1\}\gg_a X/\log X$.
The construction uses only primes for which $p-1$ has a factor
larger than $(p-1)^{1-\eta}$: all other possible index divisors then
lie in the uniform complete-splitting range, while a bad large
factor forces a small multiplicative order.  We prove the required
single-factor statistic in fixed arithmetic progressions, keeping
the character multiplier separate from the invariant marked core.
For any common width $0<\delta<1/2$ the proof gives the explicit
lower coefficient $\alpha_a(Y)[-\log(1-\delta)]-\tau(Y)$, positive
for all sufficiently large fixed $Y$.  The companion strip theorem
permits $\delta=25/258$.  No asymptotic formula or existence of a
natural density is asserted.
\end{abstract}

\section{The lower-density theorem}

An integer $a$ is \emph{admissible} if $a\ne-1$ and $a$ is not an
integer square; in particular $|a|>1$.  For $p\nmid a$, put
\[
 i_p(a)=\frac{p-1}{\ord_p(a)},\qquad
 A_a(x,2x)=\#\{x<p\le2x:p\nmid a,\ i_p(a)=1\},
\]
and write $A_a(X)$ for the corresponding count with $p\le X$.
Primes denoted by $p$ and $q$ are rational primes throughout.

We use two analytic inputs.  The first is the following common
zero-free width, for a fixed number $0<\delta<1/2$:
\begin{equation}\tag{$H_\delta$}\label{eq:H}
 L_F(s,\chi)\ne0\quad(\Re s>1-\delta),
\end{equation}
for every cyclotomic field $F$ containing $\mu_{12}$ and every
finite-order Hecke character $\chi$ of $F$; the pole of the
principal character at $1$ is permitted.  The width is independent
of $F$ and $\chi$.  This is the form of the input in
\cite[Theorem~1.2]{OAI}; the companion argument \cite{Strip}
permits $\delta=25/258$.

The second input consists of the one-sided marked Type~II and
two-sided correlation estimates of \cite[Sections~3 and 6]{PD},
together with the graph contracts used in their proofs and the
coefficient criterion in \cite[Lemma~10.2]{OAI}.  We retain their
stated bounds and invariance under removing powers of the marking
primes, including the permitted divisor growth of the invariant
core.  Appendix~\ref{sec:ap}
derives the only consequence needed here: a single long-prime
factor statistic in each fixed reduced progression.  In particular,
the argument requires neither the full Poisson--Dirichlet limit
nor a factorization law conditioned on nonabelian splitting.

Let $D$ be the signed fundamental discriminant of $\Q(\sqrt a)$.
Choose a fixed integer $Y_a\ge3$ large enough that every prime
factor of $D$ is at most $Y_a$, and that
\begin{equation}\label{eq:degree-choice}
 [\Q(\zeta_q,a^{1/q}):\Q]=q(q-1)\qquad(q>Y_a).
\end{equation}
An explicit sufficient choice is given in Lemma~\ref{lem:fields}.
For every integer $Y\ge Y_a$, define
\begin{align}
 \alpha_a(Y)
 &=\frac12\prod_{3\le q\le Y}\frac{q-2}{q-1}
 \left(1-\mu(|D|)
       \prod_{\substack{q\mid D\\q>2}}\frac1{q-2}\right),
 \label{eq:alpha}\\
 \tau(Y)&=\sum_{q>Y}\frac1{q(q-1)}.
 \label{eq:tau}
\end{align}
An empty product is one.  If $D$ is even, then $\mu(|D|)=0$.

\begin{theorem}\label{thm:main}
Assume \eqref{eq:H} and the marked analytic estimates stated above.
For every admissible integer $a$ and every fixed integer $Y\ge Y_a$,
\begin{equation}\label{eq:lower}
 \liminf_{x\to\infty}\frac{A_a(x,2x)\log x}{x}
 \ \ge\ \alpha_a(Y)[-\log(1-\delta)]-\tau(Y).
\end{equation}
The same lower bound holds for
$\liminf_{X\to\infty}A_a(X)/\pi(X)$.  Moreover,
$\alpha_a(Y)\asymp_a1/\log Y$ and $\tau(Y)\le1/Y$.
Consequently the right side of \eqref{eq:lower} is positive for
all sufficiently large fixed $Y$.
\end{theorem}

\begin{corollary}\label{cor:positive}
Under the analytic inputs of \cite{OAI,PD,Strip}, every admissible
integer $a$ is a primitive root for a set of primes of positive
lower relative density.  In particular, there are $c_a>0$ and
$x_a$ such that
\[
 A_a(x,2x)\ge c_a\frac{x}{\log x}\qquad(x\ge x_a).
\]
One may use $\delta=25/258$ in \eqref{eq:lower}.
\end{corollary}

The proof has three steps.  Fixed congruences remove all index
primes at most $Y$.  A long-prime factor statistic supplies a
positive proportion of predecessors with exactly one large factor
and a small cofactor.  Complete splitting controls the remaining
small index primes, and a product of the integers $a^j-1$ controls
the large one.  The arithmetic argument is in
Sections~\ref{sec:local}--\ref{sec:proof}; the fixed-progression
transfer of the analytic input is recorded in the appendix.

\section{Admissible congruence classes}\label{sec:local}

Put
\[
 M_Y=8\prod_{3\le q\le Y}q
\]
and let $\B_{a,Y}$ be the reduced classes $b\pmod {M_Y}$ satisfying
\begin{equation}\label{eq:classes}
 \chi_D(b)=-1,\qquad
 b\not\equiv1\pmod q\quad(3\le q\le Y),
\end{equation}
where $\chi_D$ is the primitive quadratic character of discriminant
$D$.  Since $|D|\mid M_Y$, these conditions are well defined.

\begin{lemma}\label{lem:classes}
The proportion of reduced classes modulo $M_Y$ in $\B_{a,Y}$ is
$\alpha_a(Y)$ from \eqref{eq:alpha}.  It is positive and satisfies
$\alpha_a(Y)\asymp_a1/\log Y$.  Apart from the finitely many primes
dividing $aM_Y$, every prime in these classes satisfies
\[
 \gcd\!\left(i_p(a),\prod_{q\le Y}q\right)=1.
\]
\end{lemma}

\begin{proof}
Before imposing the quadratic condition, the allowed proportion is
$P_Y=\prod_{3\le q\le Y}(q-2)/(q-1)$.  If $D$ is even, its
$2$-adic quadratic component is nontrivial on the odd classes
modulo $8$, and hence has mean zero.  The quadratic nonresidue
condition retains exactly half the allowed classes.

If $D$ is odd, then
$\chi_D(b)=\prod_{q\mid D}(b/q)$.  At a prime $q\mid D$, the
sum of $(b/q)$ over $\mathbb F_q^\times\setminus\{1\}$ equals
$-1$.  Therefore the mean of $\chi_D$ on the allowed classes is
\[
 \mu(|D|)\prod_{q\mid D}\frac1{q-2}.
\]
Using $\1_{\{\chi_D=-1\}}=(1-\chi_D)/2$ gives
\eqref{eq:alpha}.  Its bracket is strictly positive: when the
M\"obius factor is positive, $|D|$ has at least two distinct
odd prime factors, so the product in the bracket is less than
one.  The exceptional value $D=1$ would make $a$ a square.

The identity
\[
 \frac{q-2}{q-1}
 =\left(1-\frac1q\right)
  \left(1-\frac1{(q-1)^2}\right)
\]
and Mertens' theorem give $P_Y\asymp1/\log Y$, since the second
product converges to a positive limit.  This proves the claimed
size of $\alpha_a(Y)$.

For a prime in \eqref{eq:classes}, Euler's criterion gives
$a^{(p-1)/2}\equiv-1\pmod p$.  Thus $i_p(a)$ is odd.  Every
odd prime $q\le Y$ is absent from $p-1$, and hence from $i_p(a)$.
This also removes any obstruction caused by $a$ being an odd
perfect power.
\end{proof}

\section{Uniform complete splitting}\label{sec:splitting}

For a prime $q$, write
\[
 E_q=\Q(\zeta_q),\qquad K_q=E_q(a^{1/q}),\qquad
 n_q=[K_q:\Q],\qquad D_q=|\disc K_q|.
\]

\begin{lemma}\label{lem:fields}
Fix a prime $\ell\mid a$ and put $e_\ell=v_\ell(|a|)$.
For every prime $q>\max(\ell,e_\ell)$,
\begin{equation}\label{eq:fields}
 n_q=q(q-1),\qquad
 \frac{\log D_q}{n_q}\le2\log q+\log|a|.
\end{equation}
Under \eqref{eq:H}, $\zeta_{K_q}$ is zero-free in
$\Re s>1-\delta$, apart from its pole at $1$.
\end{lemma}

\begin{proof}
The prime $\ell$ is unramified in $E_q$, and the valuation of $a$
at each prime above $\ell$ is $e_\ell$, not divisible by $q$.
Thus $a$ is not a $q$th power in $E_q$.  Since $E_q$ contains
$\mu_q$, the nontrivial Kummer extension $K_q/E_q$ has degree $q$.
The relative discriminant divides the discriminant of the power
basis for $a^{1/q}$, so
\[
 D_q\le q^{q(q-2)}(q^q|a|^{q-1})^{q-1},
\]
which gives \eqref{eq:fields}.

To descend the zero-free strip, put
$F_q=\Q(\mu_{12q})$ and $\widetilde K_q=K_qF_q$.
Both $\widetilde K_q/F_q$ and $\widetilde K_q/K_q$ are abelian.
The first abelian factorization writes $\zeta_{\widetilde K_q}$
as a product of the Hecke functions covered by \eqref{eq:H}.
The second writes $\zeta_{\widetilde K_q}/\zeta_{K_q}$ as a
product of entire nonprincipal Hecke functions.  A zero of
$\zeta_{K_q}$ in the asserted region would therefore contradict
the first factorization.  This is the descent in
\cite[Lemma~9.1]{OAI}.
\end{proof}

\begin{lemma}\label{lem:split-sum}
Assume \eqref{eq:H}.  Fix $0<\eta<\delta$ and $Y\ge Y_a$.
Let $S_q(x)$ count primes in $(x,2x]$ splitting completely in
$K_q$.  Then
\begin{equation}\label{eq:split-sum}
 \limsup_{x\to\infty}\frac{\log x}{x}
       \sum_{Y<q\le(2x)^\eta}S_q(x)\le\tau(Y).
\end{equation}
\end{lemma}

\begin{proof}
For a nonnegative $f\in C_c^\infty((0,\infty))$, the smoothed
prime-ideal formula of \cite[(9.9)]{OAI} gives, for any number
field $K$ of degree $n$ and discriminant $D_K$ having the strip,
\begin{equation}\label{eq:theta}
 \Theta_{K,f}(x)
 :=\sum_{\mathfrak p}\sum_{j\ge1}
       \log N\mathfrak p\,
       f\!\left(\frac{(N\mathfrak p)^j}{x}\right)
 =x\int_0^\infty f(t)\,dt
  +O_f\!\left(x^{1-\delta}(\log|D_K|+n)\right).
\end{equation}
The constant is independent of $K$.  Indeed, the Mellin transform
of $f$ decays faster than any power on fixed strips, and the
uniform zero count is
\[
 \#\{\rho:T\le\Im\rho<T+1\}
 \ll \log|D_K|+n\log(2+|T|).
\]
The zero sum is therefore bounded by the error in
\eqref{eq:theta}; the residue at zero and the integral on
$\Re s=-1/2$ have the same bound.  This records the field
dependence needed when $q$ grows.  Uniformity in the original
theta construction is not required: the common zero-free region
and the normalized discriminant bound give uniformity here.

Given $\varepsilon>0$, choose $f\ge1$ on $[1,2]$ with
$\int f\le1+\varepsilon$.  Every completely split prime in the
dyad contributes at least $n_q\log x$ to $\Theta_{K_q,f}(x)$.
Equations~\eqref{eq:fields} and \eqref{eq:theta} imply
\[
 S_q(x)\le\frac{(1+\varepsilon)x}{q(q-1)\log x}
       +O_{a,f}\!\left(
          \frac{x^{1-\delta}\log(2q)}{\log x}\right).
\]
There are at most $(2x)^\eta$ terms, and $\eta<\delta$.
Their total error is $O_{a,f,\eta}(x^{1-\delta+\eta})
=o(x/\log x)$.  Taking the limsup and then letting
$\varepsilon\downarrow0$ proves \eqref{eq:split-sum}.
\end{proof}

For $p\nmid aq$, the condition $q\mid i_p(a)$ is equivalent to
complete splitting in $K_q$: it says that $p\equiv1\pmod q$
and that $a$ is a $q$th power in $\mathbb F_p^\times$.
All counting exceptions caused by $p\mid a$ are finite.

\section{A single large predecessor factor}\label{sec:large}

Write $P^+(n)$ for the largest prime factor of $n$, with
$P^+(1)=1$.  We need only the following one-dimensional
consequence of the prime-predecessor estimates.

\begin{proposition}\label{prop:large}
Assume the marked analytic input in Section~1.  Fix $M\ge1$,
a reduced class $b\pmod M$, and $1/2<\beta<1$.  Then
\begin{equation}\label{eq:large}
 \liminf_{x\to\infty}\frac{\log x}{x}
 \#\{x<p\le2x:p\equiv b\pmod M,
             \ P^+(p-1)>(p-1)^\beta\}
 \ge\frac{-\log\beta}{\phi(M)}.
\end{equation}
\end{proposition}

\begin{proof}
Let $f\in C_c^\infty((\beta,1))$ and
$\Phi\in C_c^\infty((1,2))$ satisfy $0\le f,\Phi\le1$.
Put
\[
 T_{x,f}(n)=\sum_{q\mid n}f\!\left(\frac{\log q}{\log x}\right),
 \qquad
 C_{x,f}=\sum_q\frac1q
                 f\!\left(\frac{\log q}{\log x}\right).
\]
The fixed-progression interior estimate proved in
Proposition~\ref{prop:interior} below, together with the prime
number theorem in this fixed progression, gives
\begin{align*}
 &\sum_{\substack{p\equiv b\pmod M}}
     \Phi((p-1)/x)T_{x,f}(p-1)\\
 &\hspace{12mm}=
 \left(\frac1{\phi(M)}\int_1^2\Phi(t)\,dt
              \int_\beta^1\frac{f(u)}u\,du+o(1)\right)
 \frac{x}{\log x}.
\end{align*}
Here $C_{x,f}\to\int f(u)\,du/u$ by the prime reciprocal sum.

For all sufficiently large $x$, the left side is bounded above
by the count in \eqref{eq:large}.  Indeed, the compact support
of $f$ has minimum $\beta+\epsilon$ for some $\epsilon>0$.
Each counted $q$ then satisfies
$q\ge x^{\beta+\epsilon}>(p-1)^\beta$; there is at most one
such factor because $\beta>1/2$.  The compact support of $\Phi$
also places $p$ strictly inside $(x,2x]$ for large $x$.
Choose $f$ and $\Phi$ increasing to one on their respective open
intervals.  Their integrals approach $-\log\beta$ and $1$,
proving the proposition.  No estimate for the factorization
boundary at exponent $1$ is needed for this lower bound.
\end{proof}

\section{Proof of the lower-density theorem}\label{sec:proof}

Fix $0<\eta<\delta$ and a permitted integer $Y$.
By Proposition~\ref{prop:large} with $\beta=1-\eta$, the primes
in the classes $\B_{a,Y}$ satisfying
\begin{equation}\label{eq:candidate}
 Q:=P^+(p-1)>(p-1)^{1-\eta}
\end{equation}
have lower dyadic coefficient at least
\begin{equation}\label{eq:candidate-mass}
 \alpha_a(Y)[-\log(1-\eta)].
\end{equation}
For such a prime, write $p-1=Qm$.  Since $\eta<1/2$,
\begin{equation}\label{eq:cofactor}
 m<(p-1)^\eta\le(2x)^\eta,
 \qquad Q>(p-1)^{1/2}.
\end{equation}
Thus $Q$ occurs only once, and every other prime divisor of
$p-1$ is at most $(2x)^\eta$.

Lemma~\ref{lem:classes} removes every index prime $q\le Y$.
If $i_p(a)>1$, a prime divisor of this index consequently falls
into exactly one of the following two ranges.

First suppose $Y<q\le(2x)^\eta$.  Complete splitting in $K_q$
and Lemma~\ref{lem:split-sum} show that the union of all these
exceptions has upper dyadic coefficient at most $\tau(Y)$.
This bound is taken over all primes, before restricting to
\eqref{eq:candidate} or to the chosen progressions.

The remaining possibility is $q=Q$.  In that case
$\ord_p(a)\le m<(2x)^\eta$.  Hence $p$ divides the nonzero
integer
\[
 \prod_{1\le j\le(2x)^\eta}|a^j-1|.
\]
Its logarithm is $O_a(x^{2\eta})$, while every counted prime
is larger than $x$.  All these exceptions together number
\begin{equation}\label{eq:large-exceptions}
 O_a\!\left(\frac{x^{2\eta}}{\log x}\right)
 =o(x/\log x),
\end{equation}
because $\eta<1/2$.  This estimate covers the large index prime
regardless of how close it is to $p$.

Subtracting both exceptional sets from \eqref{eq:candidate-mass}
proves
\[
 \liminf_{x\to\infty}\frac{A_a(x,2x)\log x}{x}
 \ge\alpha_a(Y)[-\log(1-\eta)]-\tau(Y).
\]
The parameters $Y$ and $\eta$ were fixed before $x\to\infty$.
Now let $\eta\uparrow\delta$ to obtain \eqref{eq:lower}.
By Lemma~\ref{lem:classes}, $\alpha_a(Y)\asymp_a1/\log Y$;
and the integer telescoping sum gives
\[
 \tau(Y)\le\sum_{n>Y}\frac1{n(n-1)}=\frac1Y.
\]
The coefficient is therefore positive once the fixed integer
$Y$ is sufficiently large.

Finally let $c$ denote the right side of \eqref{eq:lower}.
If $c\le0$, the asserted global lower bound is immediate.
If $c>0$, fix $J$ and sum the dyadic lower bounds over
$(X/2^j,X/2^{j-1}]$, $1\le j\le J$.  The prime number theorem
gives
\[
 \liminf_{X\to\infty}\frac{A_a(X)}{\pi(X)}
 \ge c\sum_{j=1}^J2^{-j}=c(1-2^{-J}).
\]
Letting $J\to\infty$ proves the global claim and completes
Theorem~\ref{thm:main}.
\qed

\begin{remark}
The proof uses the common strip only to ensure $\eta<\delta$
in the sum of complete-splitting errors.  Any positive common
width therefore gives positive lower density.  A wider strip
increases the explicit coefficient, but the precise value
$25/258$ is not needed for positivity.  The argument controls
all index divisors inside a deliberately restricted family of
predecessors; it does not estimate the unrestricted large-index
tail required for the full Artin asymptotic.
\end{remark}

\appendix
\section{A fixed-progression interior statistic}\label{sec:ap}

This appendix proves the fixed-modulus extension actually used in
Proposition~\ref{prop:large}.  All equation and proposition numbers
in this appendix attached to \cite{PD} refer to its
24 September 2026 version.  Its graph and marked Type~II estimates
are retained as analytic inputs.  We verify the changes caused by
the progression, including the shared-label character cancellation.

\begin{proposition}\label{prop:interior}
Fix $M\ge1$, $(b,M)=1$, a smooth compactly supported
$f:(0,1)\to\mathbb R$, and
$\Phi\in C_c^\infty((0,\infty))$.  With $T_{x,f}$ and
$C_{x,f}$ as above, the marked estimates of \cite{OAI,PD} imply
\begin{equation}\label{eq:interior}
 \sum_{\substack{p\equiv b\pmod M}}
 \Phi((p-1)/x)
       \bigl(T_{x,f}(p-1)-C_{x,f}\bigr)
 =o_{M,b,f,\Phi}(x/\log x).
\end{equation}
The modulus is fixed before $x$ tends to infinity.
\end{proposition}

We first establish the required twisted correlation.  Use the
notation of \cite[Theorem~6.1]{PD}: $F_x$ is its centered
predecessor statistic, $G$ is its invariant core, and $W_1$ is
its marked weight.  The core satisfies
\[
 G(v)=G(v^*),\qquad
 |G(v)|\le(\log x)^C\mathrm d(v^*)^C
\]
for fixed $C$, where $\mathrm d$ is the divisor function and
$v^*$ removes every prime-power factor belonging to the active
marking groups.

\begin{lemma}\label{lem:character}
For a Dirichlet character $\chi\pmod M$, extended by zero away
from units, the same estimate holds after inserting
$\chi(w+1)$ at the second endpoint:
\begin{equation}\label{eq:twisted}
 \sum_{w\ge1}\frac{\Psi(w/x)}w
 F_x(w)\chi(w+1)G(w+1)W_1(w)W_1(w+1)
 \ll_{A,M}(\log x)^{-A}.
\end{equation}
The displayed invariance and divisor-growth hypotheses are
imposed on $G$ itself.
\end{lemma}

\begin{proof}
All active marking primes eventually exceed $M$.  In the
physical pairing of \cite[Section~6.2]{PD}, the positions obey
$n'=n+D$.  Multiply the target endpoint by $\chi(n')$ and
the edge coefficient by $\overline{\chi(D)}$.  To preserve the
graph bound, expand
\[
 \overline{\chi(D)}=\sum_{h\bmod M}c_h\e(hD/M),
 \qquad \sum_h|c_h|\le\sqrt M,
 \qquad \e(t)=e^{2\pi it}.
\]
Multiplication of an edge by $\e(hD/M)$ is the unitary
conjugation $U_h^{-1}A_\zeta U_h$, where
$U_h v(n,\mathbf p)=\e(hn/M)v(n,\mathbf p)$.
Each Fourier term is therefore bounded by the original
\cite[Proposition~6.2]{PD}, applied to modulated endpoint
vectors.  Their suprema and norms do not increase; the first
endpoint remains independent of the chosen marking lists.
The residual ideal family and its comparison kernels are the
original ones.  Thus their hypotheses and logarithmic precision
are unchanged.

Now remove the shared labels as in \cite[(6.14)]{PD}:
\[
 n=D_Rw,\qquad n'=D_R(w+D_C),\qquad D=D_RD_C.
\]
Since the label products are units modulo $M$, the inserted
factors reduce exactly to
\begin{equation}\label{eq:cancellation}
 \chi(n')\overline{\chi(D)}
 =\chi(D_R)\chi(w+D_C)
     \overline{\chi(D_RD_C)}
 =\chi(w+D_C)\overline{\chi(D_C)}.
\end{equation}
No dependence on the shared product remains.  Collision and
square-exception estimates still hold, since all multipliers
have modulus at most one.

For the raw pattern, $D_C=1$, giving \eqref{eq:twisted}.
For a comparison pattern, write $D_C=D_AD_B$ as in
\cite[(6.17)]{PD}.  The factors
$\overline{\chi(D_A)}$ and $\overline{\chi(D_B)}$ are absorbed
into the arbitrary bounded free-tuple tests in (6.18).
The target sequence in (6.19) becomes
$V(v)=\chi(v)G(v)W_t^b(v)$, with the same energy bound.
That sequence is used only through its energy after the shared
labels have been removed; the original fixed divisor-moment
bounds still apply because $|\chi|\le1$.  The bilinear Fourier and local Mellin
estimates in Sections~6.3--6.6 still obtain their cancellation
from the unchanged first endpoint $F_x$.  The comparison terms
are consequently negligible to the required logarithmic
precision, proving the lemma.
\end{proof}

\begin{proof}[Proof of Proposition~\ref{prop:interior}]
The one-factor case of \cite[Section~7]{PD} extracts the
centered statistic from the preceding correlation.  We give
the changes for the fixed progression.

\emph{Initial mass and presieving.}
Impose the nonnegative condition $u+1\equiv b\pmod M$ in
\cite[Lemma~7.2]{PD}.  A presieving condition $r\mid u+1$
is incompatible with it if $(r,M)>1$.  Otherwise, for a
long-prime product $\ell$ and a divisor-model term $d'$ with
$(\ell d',r)=1$, the CRT count replacing (7.9) is
\[
 \sum_{\substack{u+1\equiv b\pmod M\\r\mid u+1,\ \ell d'\mid u}}
 \Phi(u/x)
 =\frac{x}{Mr\ell d'}\int_0^\infty\Phi(t)\,dt
   +O_{\Phi,M}(1).
\]
If $(\ell d',r)>1$, there are no solutions.  The factors
$\ell,d'$ are eventually coprime to $M$, since their prime
factors tend to infinity.  No condition on $(b-1,M)$ is needed.
The summed rounding errors and divisor-model comparisons are
unchanged apart from fixed constants.  The base mass gains
$1/M$, while deleting sieve primes dividing $M$ multiplies the
local sieve product by $M/\phi(M)$.  The net normalization is
$1/\phi(M)$; the two centered main terms still cancel.

\emph{Replacing composite prime slots.}
Expand
\[
 \1_{\{v\equiv b\pmod M\}}
 =\frac1{\phi(M)}\sum_{\chi\bmod M}
                       \overline{\chi(b)}\chi(v).
\]
In each prime-to-rough replacement, retain $\chi(m)$ in the
current prime slot $m$.  The coefficient criterion in
\cite[Lemma~10.2]{OAI}, used in the proof of the one-sided
estimate, is stable under this fixed twist.  Indeed, cancellation
against a testing character $\psi$ is replaced by cancellation
against $\chi\psi$, whose zero extension is a character of
modulus at most $\operatorname{lcm}(M,\operatorname{mod}\psi)$.
If the latter testing modulus was at most $(\log x)^B$, the
new one is at most $(\log x)^{B+1}$ for large $x$.  The
coefficient bounds, support restrictions, and allowed Mellin
twists are unchanged.  Thus the same arbitrary logarithmic
precision is available for each of the finitely many characters.

After all replacements the character product is the single
factor $\chi(v)$, with $v=u+1$.  Keep it outside the invariant
core in \cite[(7.26)]{PD}.  For fixed $v$ it is unchanged by
the allocation of active prime factors among slots.  All
failed-allocation, square, and candidate-deletion bounds use
absolute values, so their estimates are unchanged.  The final
correlation (7.29) now contains $\chi(u+1)$ and is bounded by
Lemma~\ref{lem:character}.  This proves the progression version
of the centered marked estimate (7.30).

\emph{Removing the marks.}
The prime mean-square estimate (7.33) bounds its restriction to
this progression, since the summands are nonnegative.  With
$B$ disjoint marking arrays, Cauchy--Schwarz therefore gives
the same $O_M(B^{-1/2})$ loss as in (7.34), with the fixed
number of marking groups absorbed in the constant.  First
let $x\to\infty$ for fixed array data, then let $B\to\infty$,
and finally increase the fixed sieve depth, as in Section~7.6.
The centered discrepancy tends to zero.  The argument applies
to smooth one-factor tests by the same partial summation and
Mellin separation used in (7.1).  This is \eqref{eq:interior}.
\end{proof}

\begin{thebibliography}{9}

\bibitem{OAI}
OpenAI,
\emph{Primitive roots for every admissible integer base},
4 October 2026, OpenAI Math, family 029.
\href{https://github.com/openai/math/blob/main/preprints/Primitive-roots-for-every-admissible-integer-base-October-4-2026/primitive-roots-all-integer-bases.pdf}{Repository PDF}.

\bibitem{PD}
OpenAI,
\emph{The Poisson--Dirichlet law for prime predecessors},
24 September 2026, OpenAI Math.
\href{https://github.com/openai/math/blob/main/preprints/The-Poisson-Dirichlet-Law-for-Prime-Predecessors-September-24-2026/paper.pdf}{Repository PDF}.

\bibitem{Strip}
\emph{A common zero-free strip for Hecke $L$-functions},
companion manuscript, 2026.

\end{thebibliography}
\end{document}
